\section{Absolute and relative Theorem A}

\begin{claim}{B.1}{Proved}{The component-dependent example is valid}
For all small categories put
\[
 u\in W^\dagger\quad\Longleftrightarrow\quad
 \pi _0(u)\text{ is bijective and }
 (|\pi _0(\operatorname{dom}u)|=1\Longrightarrow u\in W_\infty).
\]
This class satisfies the literal Loc~(1), Loc~(2), and absolute Loc~(3), but
not relative Loc~(3 bis). It is not closed under arrow retracts. Its strong
saturation is $W_0$.
\end{claim}
\begin{proof}
Isomorphisms belong to the class. If two of $f,g,gf$ belong, their three
component sets are in bijection. If their common cardinality is one, use
two-out-of-three for weak homotopy equivalences; otherwise the component
test is the entire test. This includes cardinality zero and arbitrary
infinite cardinalities. If $gf$ and $fg$ belong, both component maps are
bijections. In the connected case their images in the homotopy category
of spaces have invertible composites in both orders, so each factor is
invertible. This proves the moderate saturation clause.

The $W^\dagger$-aspherical categories are exactly the weakly contractible
categories. Consequently the terminal-object axiom and absolute Theorem A
follow from contraction of a category with terminal object and Quillen's
Theorem A, respectively \citep[\S1, Theorem A, p.~93]{RESQuillen1973}.

For the relative counterexample let $G=\mathbb Z/2$, let $E$ be the chaotic
groupoid on the object set $G$, and send its unique arrow $x\to y$ to
$y-x\in G$ to obtain $p:E\to BG$. The triangle is
$u:E\times BG\to E\xrightarrow p BG$. Both total categories are connected;
their fundamental groups are $G$ and $1$, so $u\notin W^\dagger$.
An object of $E/*$ is $(x,t)\in G^2$. An arrow from $(x,t)$ to $(y,t')$
exists precisely when $t'+y-x=t$. Hence $t+x$ labels two chaotic connected
components. Moreover $(E\times BG)/*=(E/*)\times BG$. The comma comparison
is a bijection on its two components and therefore belongs to
$W^\dagger$. This verifies every premise of the relative axiom and
refutes its conclusion.

The arrow $BG\to e$, which is not weak, is an arrow retract of
$(BG\amalg e)\to(e\amalg e)$. The latter belongs to $W^\dagger$; the
sections select the first components, while the retractions send the added
point to the unique object on either side. This proves failure of retract
closure. More generally, if $f:A\to B$ is bijective on components and
$A\ne\varnothing$, choose an object $a\in A$. Then $f$ is an arrow retract
of $f\amalg1_e$, using the constant maps at $a$ and $f(a)$ on the added
components. The latter belongs to $W^\dagger$. The empty case is an
isomorphism. Thus the strong saturation contains $W_0$. Conversely
$W^\dagger\subset W_0$, and $W_0$ is strongly saturated as the inverse
image of isomorphisms under $\pi _0$. This proves equality.
\end{proof}

\begin{claim}{B.2}{Proved}{An elementary quotient of the homotopy category}
Let $\mathscr K$ be the homotopy category of CW homotopy types, and fix
$K=S^1$. Let $I(X,Y)$ be the set of homotopy classes $X\to Y$ factoring
through $K$. There is a category $\mathscr Q$ with the same objects and
\[
 \mathscr Q(X,Y)=\bigl(\mathscr K(X,Y)\amalg\{0\}\bigr)/
       \bigl(I(X,Y)\sim0\bigr).
\]
For the quotient functor $q:\mathscr K\to\mathscr Q$, the morphism $q(f)$
is invertible exactly when $f$ is invertible in $\mathscr K$, or both its
source and target are homotopy retracts of $S^1$. Those retracts have
homotopy type $e$ or $S^1$.
\end{claim}
\begin{proof}
The family $I$ is a two-sided ideal: precomposition or postcomposition
preserves a factorization through $K$. Declare the adjoined zeros
absorbing. Composition therefore descends, is associative, and preserves
the quotient identities. We adjoin a zero even when the original hom-set
is empty. In particular the empty space causes no exception.

The equation $q(1_X)=0$ holds precisely when $1_X$ factors through $K$,
which is exactly the assertion that $X$ is a homotopy retract of $K$.
Such an object is a zero object of $\mathscr Q$: every morphism to or from
it factors through $K$, and hence every corresponding hom-set is a
singleton. Conversely, suppose $q(f):X\to Y$ is invertible and neither
identity is zero. An inverse cannot be the zero morphism. Represent it by
$g:Y\to X$ in $\mathscr K$. Every nonzero equivalence class in the
quotient is a singleton, so $q(gf)=q(1_X)$ and $q(fg)=q(1_Y)$ imply the
actual equalities $gf=1_X$ and $fg=1_Y$ in $\mathscr K$. If exactly one
identity is zero, invertibility is impossible; if both are zero it is
automatic. These observations prove the criterion.

If $X$ is a homotopy retract of $S^1$, it is nonempty and connected, its
higher homotopy groups vanish, and $\pi _1X$ is a retract of $\mathbb Z$.
An idempotent endomorphism of $\mathbb Z$ is multiplication by $0$ or $1$.
Thus $\pi _1X$ is $0$ or $\mathbb Z$. The CW Whitehead theorem identifies
$X$ with $e$ or $S^1$. Both types are indeed retracts. The empty space is
not one, since there is no map $S^1\to\varnothing$.
\end{proof}

\begin{claim}{B.3}{Refuted}{Retracts and strong saturation do not repair the absolute axiom}
Neither Loc~(1 bis), Loc~(2), Loc~(3) nor Loc~(1 ter), Loc~(2), Loc~(3)
implies Loc~(3 bis). A counterexample on all small categories is
\[
 W_{\mathrm{circ}}=W_\infty\ \cup\
 \{u:A\to B\mid BA\simeq S^1\simeq BB,
          \ H_1(Bu;\mathbb Q)\text{ is invertible}\}.
\]
In the second alternative, the induced integer degree is any nonzero
integer. The $W_{\mathrm{circ}}$-aspherical categories are exactly the
weakly contractible categories.
\end{claim}
\begin{proof}
Consider the ordinary functor
\[
 \Phi:\mathscr K\longrightarrow\mathscr Q\times\operatorname{Vect}_{\mathbb Q},
 \qquad X\longmapsto(qX,H_1(X;\mathbb Q)).
\]
By B.2, the inverse image of isomorphisms consists of homotopy
equivalences, together with maps between retracts of $S^1$ which are
isomorphisms on rational first homology. Among those retracts, a map from
the contractible type to the circle type, or conversely, fails the latter
test; between contractible types it is already a homotopy equivalence;
between circles the test is exactly nonzero degree. Classifying spaces of
categories are CW complexes, so the inverse image of isomorphisms under
$\Phi B$ is exactly the displayed class.

An inverse image of isomorphisms under any functor is strongly saturated:
the functor factors through localization at that inverse image, and any
arrow inverted by the localization must consequently have invertible
image. It is therefore closed under retracts and satisfies all clauses of
Loc~(1 ter), including the literal moderate saturation. A single
universe enlargement suffices for the homotopy category and its quotient;
the formula defining $W_{\mathrm{circ}}$ itself quantifies over the
original small categories.

Since the target of $A\to e$ is not circle type, this arrow belongs to
$W_{\mathrm{circ}}$ exactly when $BA$ is weakly contractible. Loc~(2)
follows. If every $A/b$ is $W_{\mathrm{circ}}$-aspherical, Quillen's
Theorem A makes $A\to B$ a weak homotopy equivalence, hence a member of
$W_{\mathrm{circ}}$. This verifies absolute Loc~(3) on all small
categories. The finite relative counterexample is B.4.
\end{proof}

\begin{claim}{B.4}{Proved}{A finite certificate for the failure of relative Theorem A}
The class in B.3 violates relative Theorem A on a functor from a ten-object
poset to a six-object poset, over a three-object poset. The failure is
detected integrally by multiplication by $2$ on $H_2$.
\end{claim}
\begin{proof}
For even $m\ge4$ let $C_m$ have objects $0,\ldots,m-1$, with the
nonidentity comparisons
\[
 2i<2i-1,\qquad 2i<2i+1\pmod m.
\]
There are no composable nonidentity arrows. Its nerve is the polygon
with $m$ edges. The order map $v:C_8\to C_4$, $i\mapsto i\bmod4$,
wraps this polygon twice around the target, and belongs to
$W_{\mathrm{circ}}$. We can check its degree without choosing a
topological parametrization: the cycle
\[
 z_m=\sum_{i\text{ even}}\bigl([i,i+1]-[i,i-1]\bigr)
\]
generates $H_1(NC_m;\mathbb Z)$, and $v_\#z_8=2z_4$.

Let $\widehat C_m$ be obtained by adjoining two incomparable objects
$t_+$ and $t_-$ above every object of $C_m$. Its nerve is the union of
two simplicial cones on $NC_m$, hence is a triangulated $S^2$. Let
$S=\{o<+, o<-\}$ and map $C_m$ to $o$ and $t_\pm$ to $\pm$.
The order map $\widehat v:\widehat C_8\to\widehat C_4$ extends $v$
and preserves the two added objects. At $o$, its overcomma comparison
is $v$. At each of $+$ and $-$, the overcomma category is the corresponding
single cone, with terminal vertex $t_\pm$. Those two comparisons are
weak homotopy equivalences. Thus every relative premise holds.

The top cycle obtained by coning $z_m$ to $t_+$ and subtracting its cone
to $t_-$ is primitive and generates $H_2(N\widehat C_m;\mathbb Z)$.
The induced map sends it to twice the target generator. Thus
$\widehat v$ is not a weak equivalence, and its endpoints are sphere
types rather than circle types; it is not in $W_{\mathrm{circ}}$.

The independent program \texttt{code/BC/verify\_counterexamples.py}
constructs these orders and chain maps, checks $d^2=0$, computes all
boundary Smith normal forms through degree four, verifies the two
displayed degree equations, and computes fundamental-group presentations
from maximal trees and triangle relations. Its Tietze eliminations leave
one free generator for each crown and none for either suspension or cone.
The output contains the relations and each elimination. The simplicial
description already proves the asserted homotopy types; the program is
an independent exact certificate, not a homology-only inference of them.
\end{proof}

\begin{claim}{B.5}{Known}{Literal homotopy invariance and Rappel 1.2}
Under literal Loc~(1), Loc~(2), and absolute Loc~(3), naturally homotopic
functors have the same image in $\operatorname{Hot}_W$ and either both or
neither belong to $W$. Every category-homotopy equivalence belongs to
$W$. In particular both functors in an adjunction belong to $W$, and a
category with initial object is $W$-aspherical.
\end{claim}
\begin{proof}
For every $A$, the overcomma categories of
$p:A\times[1]\to A$ are $(A/a)\times[1]$ and have terminal objects.
Absolute Theorem A therefore gives $p\in W$. Both endpoint sections
$i_0,i_1$ lie in $W$ by two-out-of-three, and their images in the
localization equal the same inverse of $p$. A natural transformation
$f\Rightarrow g$ is a functor $H:A\times[1]\to B$ with
$Hi_0=f$, $Hi_1=g$. These identities prove both assertions about $f,g$;
a zigzag of transformations gives the same conclusion.
If $f,g$ are category-homotopy inverse functors, $gf$ and $fg$ belong
to $W$ by comparison with the identities. The moderate clause gives
$f,g\in W$. An adjunction supplies the required homotopies by its unit
and counit. For an initial object the natural transformation from its
constant functor to the identity contracts the category.
This is the elementary argument underlying the homotopy results in
\citet[\S1.1]{RESMalts2020}; no relative Theorem A is used.
\end{proof}

\begin{claim}{B.6}{Known}{Rappel 1 holds for the literal axioms}
Literal Loc~(1), Loc~(2), Loc~(3) imply $W=W^\circ$.
\end{claim}
\begin{proof}
Let $\operatorname{Tw}(A)$ have objects arrows $f:x\to y$ and
morphisms $f\to f'$ pairs $u:x'\to x$, $v:y\to y'$ with $f'=vfu$.
Its source and target projections are
\[
 A^{\mathrm{op}}\xleftarrow{s_A}\operatorname{Tw}(A)
       \xrightarrow{t_A}A.
\]
The category $t_A/a$ consists of composable pairs $x\xrightarrow f y
\xrightarrow t a$. Composition gives a functor to $(A/a)^{\mathrm{op}}$;
its section sends $h:x\to a$ to $(h,1_a)$. The arrows
$(1_x,t)$ give a natural transformation from the identity to section
followed by composition. Hence this is a category-homotopy equivalence.
The target has initial object $1_a$, so B.5 and two-out-of-three make
$t_A/a$ $W$-aspherical.

Similarly, an object of $s_A/a$, with the latter comma taken in
$A^{\mathrm{op}}$, is $a\xrightarrow h x\xrightarrow f y$.
Composition maps it to $a\backslash A$, with section
$a\xrightarrow{1_a}a\to y$. The arrows $(h,1_y)$ again give a
natural transformation from the identity to section followed by
composition. The category $a\backslash A$ has initial object.
Absolute Theorem A gives $s_A,t_A\in W$. The construction is functorial
in $A$. For $u:A\to B$, the squares with $s,t$ and
$\operatorname{Tw}(u)$ commute. Two applications of two-out-of-three
therefore give $u\in W\Longleftrightarrow\operatorname{Tw}(u)\in W
\Longleftrightarrow u^{\mathrm{op}}\in W$.
This is the twisted-arrow proof of
\citet[Proposition~1.1.21 in the March 2020 version]{RESMalts2020};
the same construction occurs in \citet[\S1, p.~94]{RESQuillen1973}.
\end{proof}

\begin{claim}{B.7}{Proved}{The two equivalences in Proposition 1.3}
Both equivalences in XVI Proposition~1.3 hold when its category
$\mathcal M$ is read as $\operatorname{Cat}$. The weak equivalence uses
absolute Theorem A; the strong equivalence uses relative Theorem A.
\end{claim}
\begin{proof}
For a Grothendieck fibration $p:X\to S$, inclusion of the strict fibre
$X_s$ in $s\backslash X$ has a right adjoint, given by cartesian
transport. B.5 makes that inclusion weak. The equivalence of absolute
Theorem A with its opposite form follows from B.6. Consequently
aspherical fibres imply $p\in W$.

The split replacement $\Phi_S(X)$ has objects $(s,x,\alpha:s\to p(x))$.
A morphism is a pair $(h:s\to t,k:x\to y)$ with
$\beta h=p(k)\alpha$. It is the split Grothendieck fibration of
$s\mapsto s\backslash X$. The comparison
$i_X(x)=(p(x),x,1)$ has a left adjoint $(s,x,\alpha)\mapsto x$,
and therefore belongs to $W$. This proves the forward weak implication.
For a map between two fibrations, the fibre inclusions give a commutative
square comparing fibre maps with the maps on $s\backslash X$.
Two-out-of-three and the opposite form of relative Theorem A prove the
forward strong implication. No preservation of chosen cartesian arrows
is needed for this last square.

Conversely, assume the segment axiom and the split-fibration clauses.
The projection $A\times[1]\to A$ is a split fibration with aspherical
fibres, hence lies in $W$. The proof of B.5 now applies verbatim, without
assuming Theorem A, and proves the terminal- and initial-object axioms.
If all $s\backslash X$ are aspherical, the split replacement has
aspherical fibres, so $\Phi_S(X)\to S$ and $i_X$ are weak. Thus
the opposite absolute axiom holds. Apply B.6 to $W^\circ$, which now
satisfies the absolute axioms: this gives $W=W^\circ$ and hence the
absolute axiom for $W$ itself. This step removes the apparent circularity
in the manuscript's backward proof. For the strong clause use
$\Phi_S(f)$ and the natural weak comparisons $i_X,i_Y$; the same fibrewise
argument gives the opposite relative axiom, hence the relative axiom.
\end{proof}

\begin{claim}{B.8}{Refuted}{Products are not consequences of the literal definition}
Even under Loc~(1 ter), Loc~(2), Loc~(3), finite products of weak arrows
need not be weak. Thus XVI Rappel~2.3, and the finite-product assertion
of Proposition~2.4 under the printed definition alone, require correction.
\end{claim}
\begin{proof}
For $W_{\mathrm{circ}}$, take degree two on the circle and multiply it
by the identity of a second circle. The endpoints have torus type; the
map is not a weak equivalence, since its map on $\pi _1$ has image of
index two. It therefore does not belong to $W_{\mathrm{circ}}$.
If localization preserved binary products, the product of an inverted
arrow and an identity would be inverted; strong saturation would put
this map in $W_{\mathrm{circ}}$, a contradiction. The same example
adjoined to an isolated point proves failure of finite-coproduct
closure under the literal definition. The editorial qualification at
Rappel~2.3 correctly supplies relative Loc~(3 bis).
\end{proof}

\begin{claim}{B.9}{Known}{The exact literature distinction}
Every class satisfying the literal package is a weak fundamental
localizer in the sense of Cisinski~6.1.12 and Maltsiniotis~1.1.2.
Their saturation terminology uses a split-idempotent clause implied
by the literal moderate clause; their Theorem A clause is exactly
absolute. The sources inspected do not supply the counterexample of B.3.
\end{claim}
\begin{proof}
We checked \citet[\S6.1.12]{RESCisinski2006}: its clauses are weak
saturation, the terminal-object axiom, and the assertion that every
aspherical functor is weak. Here weak saturation means FS1--FS3 of
Cisinski~1.4.12: identities, two-out-of-three, and
$ri=1,\ ir\in W\Rightarrow i\in W$. The literal moderate clause
implies FS3; we do not identify these two formulations in an arbitrary
category. In particular FS3 is not closure under arbitrary arrow
retracts. The March 2020 version of
\citet[\S1.1.2]{RESMalts2020} writes the third premise as
$A/b\to B/b\in W$; because $B/b$ has terminal object this is precisely
the absolute axiom. There is no additional product or descent clause.
\citet[Remark~1.3]{RESMalts2012} states:
``On ne conna\^it pas d'exemple de localisateur fondamental faible qui ne
soit pas un localisateur fondamental.'' The original author's
MathOverflow answer of December 2019 states the same absence of a known
example and points to Cisinski~6.1.11, 6.1.17, and 6.1.18
\citep{RESCisinski2019MO}. These are historical reports, not a claim about
everything published through the present date. We searched the exact
French and English terms, their conjunctions with absolute and relative
Theorem A, and the cited primary texts; no priority is claimed for B.3.
\end{proof}

\begin{claim}{B.10}{Open}{The uncompleted statement-by-statement historical audit}
The complete audit of every assertion in XVI ms.~1--27 and XV~\S9,
including all implications among the additional Mayer--Vietoris axioms,
is not completed here. The complete list of settled and unresolved
named assertions is \texttt{sections/BC\_historical\_audit.csv}.
B.11--B.14 below settle the formal decomposition statements, finite
localization statements, the two discrete-set adjunctions, and the
elementary cube implications. The remaining unresolved assertions
include the infinite-product generality of Lemma~2.5 after adding
identities, the variants with the full decomposition hypothesis in
XV~9.7 and 9.11, and the full Mayer--Vietoris implications in
XVI~3.2--3.7. B.17 refutes several of the other nontriviality
equivalences under the literal absolute axioms. Replacing absolute
Loc~(3) in D\'efinition~1.1 by relative Loc~(3 bis) repairs the product
argument in B.8; the separate correction in B.11 is still needed.
Optional axioms Loc~(4)--(8) retain their stated hypotheses.
\end{claim}
The quotient example proves the missing descent condition very concretely:
$\Phi$ in B.3 inverts degree two on $S^1$, but does not invert its
suspension. Thus an argument through homotopy colimits cannot follow
from strong saturation alone. The group-dependent universal-cover
homology proposal is unnecessary for the negative answer. We have not
asserted that every possible version of that proposal has been classified.

\begin{claim}{B.11}{Refuted}{The parenthesis in XVI Corollary 2.2 needs a separate correction}
Closure under finite coproducts and filtered colimits does not imply
the full decomposition axiom Loc~(4), even for a strongly saturated
relative fundamental localizer. The parenthesis in Corollary~2.2 must
say Loc~(4 a), or must separately assume Loc~(4 b).
\end{claim}
\begin{proof}
Let $\epsilon:\operatorname{Cat}\to[1]$ send the empty category to $0$
and every nonempty category to $1$, and put
\[
 W_{\mathrm{ne}}=\epsilon^{-1}(\operatorname{Iso}[1]).
\]
This is the coarse localizer of
\citet[Example~1.1.30]{RESMalts2020}.
An existing functor is weak precisely when its source and target
are both nonempty or both empty. This inverse-image description proves
strong saturation. The terminal-object axiom is immediate. For relative
Theorem A, suppose $u:A\to B$ lies over $S$ and every $u/s$ is weak.
If $B$ is empty, so is $A$. If $B$ is nonempty, choose $b\in B$ and
put $s=p(b)$. The comma category $B/s$ contains $(b,1_s)$, hence $A/s$
is nonempty and therefore $A$ is nonempty. In either case $u$ is weak.
This verifies the axiom on all small categories.

A coproduct is nonempty exactly when at least one summand is nonempty,
so arbitrary coproducts preserve this weak class. A filtered colimit of
categories is nonempty exactly when one of its terms is nonempty; this
follows already on object sets. Thus filtered colimits preserve the
class too. Nevertheless the inclusion $e\to e\amalg e$ is weak,
whereas its restriction to the missed summand is $\varnothing\to e$,
which is not. This is a finite counterexample to Loc~(4 b). Its category
tables, the inclusion, and its strict pullback to the other summand are
checked in \texttt{verify\_counterexamples.py}.
\end{proof}

\begin{claim}{B.12}{Proved}{The formal decomposition and finite localization assertions}
XVI Proposition~2.1, the principal equivalence of Corollary~2.2,
the finite-product and finite-coproduct versions of Lemma~2.5 when
identities belong to $W$, and Lemma~2.6 hold without relative Theorem A.
The identity hypothesis is automatic for every localizer considered
here; B.15 treats the lemma's more general literal wording.
The formal equivalences in XV Proposition~9.4,
Corollary~9.5, and Corollary~9.6 also hold, except that this statement
does not certify 9.4's parenthetical deduction from earlier axioms.
XVI Corollary~2.9 at ms.~13 follows from summand restriction alone.
\end{claim}
\begin{proof}
Write $\Sigma$ for closure of $W$ under the indicated coproducts and
$D$ for restriction of a weak arrow to a summand of its codomain.
Every arrow over a coproduct of categories is the coproduct of its
restrictions. Hence the assertion that a coproduct is weak exactly
when each summand is weak is equivalent to $\Sigma+D$. Binary closure
is equivalent to preservation under $f\mapsto f\amalg1_Z$: one
direction uses the identity arrow; the other factors
$f\amalg g=(1\amalg g)(f\amalg1)$. Induction gives finite closure.
A coproduct of an arbitrary family is the filtered colimit of its
finite subcoproducts. Thus filtered-colimit closure upgrades finite
$\Sigma$ to arbitrary $\Sigma$, but does not provide $D$, as B.11 shows.
Restriction to any one summand follows already by writing a coproduct
as that summand plus the sum of the others, even for an infinite family.
Under finite closure, restriction to the finitely many components of
a summand implies restriction to the summand; filtered-colimit closure
gives the same conclusion with infinitely many components.
These observations prove the stated formal decomposition equivalences.

For the component test, let $f:X\to Y$ be weak and restrict to each
connected component $Y_j$. By $D$ its restriction is weak. If weak
arrows into connected categories have connected, nonempty sources,
each inverse image $X_j$ is connected and nonempty. Consequently
$\pi_0f$ is a bijection. The converse is immediate. This proves both
versions of the stated connectivity criterion.

If $F:\mathcal A\rightleftarrows\mathcal B:G$ is an adjunction and
both functors preserve the specified weak arrows, its unit and counit
descend to the localizations. Their naturality on inverse weak arrows
follows by multiplying the original naturality squares by inverses.
The triangle identities remain identities, so the descended functors
are adjoint. This proves XVI~2.6.

For two relative categories, localization of their product at the
componentwise weak arrows is the product of their localizations.
Indeed a functor inverting those arrows inverts weak arrows in each
coordinate separately. Extend first in one coordinate and then in the
other; the commutation identities persist when an arrow is inverted,
which constructs the unique extension to the product localization.
The same argument applies to natural transformations. Now descend the
adjunction between the diagonal and binary product, assuming product
closure of $W$. It proves the binary case of XVI~2.5; induction proves
the finite case. A terminal object survives any localization: the
unique ordinary maps to it define a natural transformation from the
identity to its constant functor, whose naturality extends across
inverted arrows; at the terminal object its component is the identity.
Hence every localized map into it is that prescribed component.
Taking opposites gives the
finite-coproduct and initial-object assertions. These arguments also
prove the finite-coproduct clause of Proposition~2.4 when its explicit
closure hypothesis holds. No identification of an infinite product of
localizations with localization of an infinite product has been used.
\end{proof}

\begin{claim}{B.13}{Known}{The two adjunctions with discrete sets need only the literal axioms}
Under literal Loc~(1), Loc~(2), Loc~(3) and $W\subseteq W_0$, the
functor $\pi_0:\operatorname{Hot}_W\to\operatorname{Set}$ is both
left and right adjoint to the discrete-set functor. Thus the second
Proposition~2.8, the second Corollary~2.9, and Corollaries~2.11--2.13
of XVI hold as printed, for arbitrary small sets. Lemma~2.14 also
holds in its stated abstract generality.
\end{claim}
\begin{proof}
Cisinski's minimality theorem gives $W_\infty\subseteq W$
\citep[Theorem~6.1.18]{RESCisinski2006}. Its saturation hypothesis
is precisely FS1--FS3 of \S1.4.12; FS3 says that $ri=1$ and
$ir\in W$ imply $i\in W$, and is implied by the literal moderate
clause. Thus no arrow-retract axiom has been inserted here.
Localize first at $W_\infty$, obtaining the homotopy category of
spaces \citep[Example~4.1.29]{RESCisinski2006}, and then at the
images of $W$. For every discrete small
set $I$ and every space $X$ one has
\[
 [I,X]=\operatorname{Set}(I,\pi_0X),\qquad
 [X,I]=\operatorname{Set}(\pi_0X,I).
\]
Both represented functors invert $W$, since $W\subseteq W_0$.

We spell out why a represented functor which already inverts $W$
keeps its representing hom-sets after localization. If
$\mathcal A(c,w)$ is bijective for every $w\in W$, a zigzag
starting at $c$ can be shortened inductively: replace
$c\xrightarrow h y\xleftarrow w z$ by the unique $k:c\to z$
with $wk=h$. Thus every localized arrow from $c$ is represented
in $\mathcal A$. Injectivity follows by extending
$\mathcal A(c,-)$ to the localization and evaluating a localized
arrow at $1_c$. The dual argument treats arrows into $c$.
Apply this twice to the discrete object $I$ in the homotopy category
of spaces. The two displayed formulas survive the second localization
and are precisely the two adjunctions. In particular the discrete-set
functor is fully faithful, $\operatorname{Hot}_W(e,X)=\pi_0X$,
every arrow from $e$ is represented by an object of $X$, and there
is no arrow from a nonempty $X$ to the empty category.

For the abstract Lemma~2.14, the representable functor
$\mathcal M(\varnothing,-)$ is constantly a singleton and inverts
every weak arrow. The preceding representable argument shows that
the initial object remains initial. If no weak arrow has empty source and nonempty
target, a zigzag ending at the strict initial object must remain at
that object, by backwards induction: an ordinary arrow into it has
initial source, and an inverse weak arrow into it has initial source
by the hypothesis on $W$. Hence it remains strict. To prove the
asserted surjectivity from the terminal object $e$, shorten each
segment $e\xrightarrow{u_1}X_1\xleftarrow{u_2}X_2$, $u_2\in W$,
using the stipulated $s:e\to X_2$ with $u_2s=u_1$ in the
localization. Repeating removes every backward arrow. The result
is represented by one ordinary arrow $e\to X$, exactly as claimed.
\end{proof}

\begin{claim}{B.14}{Proved}{The elementary cube implications are formal}
The three implications of XVI Proposition~3.1 hold under Loc~(1).
They require neither absolute nor relative Theorem A.
\end{claim}
\begin{proof}
Start with a Mayer--Vietoris square having maps
$i_1:X_0\to X_1$ and $j_1:X_2\to X_3$ on opposite sides.
Map the degenerate Mayer--Vietoris square on $X_0\to X_2$
to this square. Its four comparison maps are
$(1_{X_0},i_1,1_{X_2},j_1)$. Loc~(6 a) gives
$i_1\in W\Rightarrow j_1\in W$, and Loc~(6 c) gives
the reverse implication. Next map the given square to the
degenerate square on $X_1\to X_3$. The four maps are
$(i_1,1_{X_1},j_1,1_{X_3})$. Loc~(6 b) gives the reverse
implication, and Loc~(6 c) gives the forward implication.
Every degenerate square is still a union-of-two-sieves square,
with one sieve equal to the whole category. This proves the
three implications, including both consequences of Loc~(6 c).
The second cube's repeated label $f_1$ in the transcription must
read $f_2=j_1$ at its third vertex.
\end{proof}

\begin{claim}{B.15}{Refuted}{The arbitrary-arrow wording of XVI Lemma 2.5 is too broad}
If $W$ is literally an arbitrary class of arrows, as specified in
XVI Lemma~2.5, even its closure under arbitrary products does not
imply that localization preserves binary products. The finite version becomes
true after requiring identities to belong to $W$, as in B.12.
\end{claim}
\begin{proof}
Let $\mathcal M$ be the four-object diamond poset
$0<a<1$, $0<b<1$, with $a,b$ incomparable, and put
$W=\{f:a\to1,1_1\}$. Arbitrary products are meets. A product of
arrows in $W$ is $f$ if any factor is $f$, and is $1_1$ otherwise,
including the nullary case. Thus every literal product-closure
hypothesis holds. But $f\times1_b$ is
$g:0\to b$. The order map to the three-object chain $0<1<2$ which
sends $0$ to $0$, $b$ to $1$, and $a,1$ to $2$ sends $f$ to an
identity and $g$ to the noninvertible arrow $0\to1$.
Consequently $g$ is not inverted by localization at $W$.
If that localization preserved binary products, it would invert
$f\times1_b$, since it inverts $f$. This contradiction proves the
claim. The program enumerates the diamond's complete composition
table and its sixteen binary meets, checks the detecting functor,
and checks the four binary products of weak arrows and the
displayed product with $1_b$.
\end{proof}

\begin{claim}{B.16}{Proved}{A weak localizer with near-bijective infinite component maps}
Call a map of sets $f:I\to J$ a near bijection if there are finite
subsets $I_0\subset I$ and $J_0\subset J$ such that its restriction
$I\setminus I_0\to J\setminus J_0$ is a bijection. Define, on all
small categories, $W_{\mathrm{nb}}$ by the following two alternatives:
\[
\begin{split}
u:A\to B\in W_{\mathrm{nb}}\quad\Longleftrightarrow\quad&
\bigl(\pi_0A,\pi_0B\text{ finite and }u\in W_\infty\bigr)\\
&\text{or }\bigl(\pi_0A,\pi_0B\text{ infinite and }
                 \pi_0u\text{ a near bijection}\bigr).
\end{split}
\]
This class satisfies literal Loc~(1), Loc~(2), absolute Loc~(3),
and closure under finite coproducts. Its aspherical categories are
exactly the weakly contractible categories. There are distinct
parallel arrows in its localization even though it contains maps
which are not bijective on components.
\end{claim}
\begin{proof}
We first construct a detecting category. It has the small sets as
objects. Adjoin an absorbing zero to every hom-set of
$\operatorname{Set}$, identify every finite-image function with zero,
and identify two functions $f,g:I\to J$ when
\[
 \{f(i),g(i)\mid i\in I,\ f(i)\ne g(i)\}
 \quad\hbox{is finite}.
\tag{B.16.1}
\]
Condition (B.16.1) is an equivalence relation. For transitivity, the
unequal values of $f,h$ lie in the union of the unequal-value sets
for $f,g$ and $g,h$. Precomposition only restricts the set of unequal
values, and postcomposition sends it to the image of a finite set.
Finite-image functions form a two-sided ideal. The zero identifications
are therefore compatible with this congruence and with composition.
We obtain a category $\mathscr R$ and a functor
$q:\operatorname{Set}\to\mathscr R$.

The identity of a finite set is zero, so every finite set is a zero
object of $\mathscr R$, including the empty set. The identity of an
infinite set is nonzero. Thus finite and infinite sets cannot become
isomorphic. For infinite $I,J$, the arrow $q(f)$ is an isomorphism
if and only if $f$ is a near bijection. To prove necessity, a nonzero
inverse is represented by a function $g:J\to I$, and (B.16.1) gives
$gf=1_I$ and $fg=1_J$ outside finite subsets. Consequently $f$ is
injective outside a finite set, every fibre of $f$ is finite, and
its image is cofinite. If $I_0$ is a finite exceptional set for
injectivity, remove
$J_0=(J\setminus f(I))\cup f(I_0)$ and then its finite inverse image.
The remaining restriction is bijective. Conversely, invert a near
bijection off its two finite exceptional sets and extend the inverse
arbitrarily on the exceptional codomain. The two composites differ
from the identities on finite subsets, so they equal the identities
in $\mathscr R$.

We verify the saturation clauses without any implicit finiteness
assumption. If two of $f,g,gf$ lie in $W_{\mathrm{nb}}$, all three
categories have finite component sets, or all three have infinite
component sets. In the first case use two-out-of-three for
$W_\infty$; in the second use the isomorphism criterion for $q\pi_0$.
If $gf$ and $fg$ lie in $W_{\mathrm{nb}}$, their images under
$q\pi_0$ are isomorphisms, so the images of $f,g$ are isomorphisms.
Their component sets therefore are either both finite or both
infinite. In the finite case the moderate clause for weak homotopy
equivalences applies; in the infinite case the near-bijection
criterion applies. Identities and isomorphisms satisfy the definition.
This proves every clause of literal Loc~(1), including empty categories.

The class contains $W_\infty$, and an arrow to the point belongs to
it precisely when it is a weak homotopy equivalence. The terminal
axiom and absolute Theorem A follow, the latter from Quillen's
Theorem A. For a finite family of weak arrows, either all its
categories have finitely many components and ordinary weak
equivalences give the conclusion, or a summand has infinitely many
components. In that case the coproduct component map is a near
bijection: take the union of the finitely many exceptional finite
sets; the finite-component summands have bijective component maps.
The empty coproduct is the identity of the empty category.

Regard $\mathbb N$ as a discrete category. The shift
$s(n)=n+1$ is in $W_{\mathrm{nb}}$. On the other hand, the functor
$q\pi_0:\operatorname{Cat}\to\mathscr R$ inverts every member of
$W_{\mathrm{nb}}$, and distinguishes $1_{\mathbb N}$ from the
constant-zero endofunctor: the former has nonzero identity image,
whereas the latter has finite image and is zero. It factors through
the localization and proves that these parallel arrows remain
distinct. All categories, arrows, and compositions in this
counterexample are explicitly discrete; the infinity of
$\mathbb N$ is essential to the proof and is not inferred from a
finite simulation. The program's exact affine-tail representation
checks the shift, its near inverse, and their finite exceptional
sets as functions on all natural numbers.
\end{proof}

\begin{claim}{B.17}{Refuted}{Several nontriviality equivalences fail for the printed definition}
Under literal Loc~(1), Loc~(2), Loc~(3), XVI Proposition~2.7
and XV Proposition~9.9, Lemma~9.10, Proposition~9.12, and
Corollary~9.13 are false. The extra finite-coproduct hypothesis
of XV~9.9 does not repair it. Relative Loc~(3 bis) repairs these
assertions.
\end{claim}
\begin{proof}
Use $W_{\mathrm{nb}}$. Its shift $s:\mathbb N\to\mathbb N$
fails the component-bijection axiom. Nevertheless its localization
and its full subcategory of nonempty categories have distinct
parallel endomorphisms of $\mathbb N$, by B.16. Neither category
is equivalent to a point, and the full localization is not
equivalent to $[1]$. This directly contradicts XVI~2.7 and
XV~9.9; the latter's finite-coproduct hypothesis was proved in B.16.
The map $t:\mathbb N\to\mathbb N$ defined by
$t(0)=t(1)=0$ and $t(n)=n-1$ for $n\ge2$ is a near bijection.
It is inverted but is not injective on components, disproving
Lemma~9.10. The distinct arrows $1_{\mathbb N}$ and the constant
endofunctor directly disprove Corollary~9.13.

For Proposition~9.12, its conditions (c), (c'), and (c'') hold,
whereas (a) fails. Its condition (d) also holds: infinite sets of
different cardinalities cannot be isomorphic in $\mathscr R$,
because a near bijection preserves infinite cardinality. The
detecting functor therefore separates arbitrarily many infinite
cardinalities already in $\operatorname{Hot}_{W_{\mathrm{nb}}}$.
Its condition (b) fails explicitly. Let $g_0,g_1:\mathbb N\to\{0,1\}$
be respectively the constant-zero map and the map with value $1$
only at $0$. Then $g_0s=g_1s$, and inversion of $s$ makes
$g_0=g_1$ in the localization, although their component maps differ.

We also give the corrected implication. Relative Theorem A
implies $f\times1_X\in W$ for $f\in W$: apply it over $X$,
whose comma comparisons are $f\times1_{X/x}$.
The projections from the source and target of the latter comparisons
are category-homotopy equivalences, since $X/x$ has terminal object.
Thus those comparisons belong to $W$ by B.5 and two-out-of-three.
B.12 now proves that localization preserves finite products.
If a weak $f:I\to J$ identifies two distinct components, select
$i_0,i_1$ in them with images in one component of $J$. Any two
functors $u_0,u_1:X\to Y$ extend to a functor $I\times X\to Y$
using the decomposition of $I$ separating $i_0,i_1$. Inverting
$f\times1_X$ transfers this map to the localization of $J\times X$.
The two sections $X\to J\times X$ at $f(i_0),f(i_1)$ are
category-homotopic, hence $u_0=u_1$ in the localization.
If instead a weak $f$ misses a component, the two maps
$J\to\{0,1\}$ which agree on its image and differ on that
component become equal. Evaluating at a point of the missed
component makes the two point inclusions into $\{0,1\}$ equal.
Product preservation then makes the two inclusions
$X\to X\amalg X$ equal for every $X$, and again every pair
$u_0,u_1$ becomes equal.

Once all ordinary parallel functors are identified, all nonempty
categories become canonically isomorphic by their constant functors.
The category obtained by identifying all parallel ordinary functors
has one initial empty object and one isomorphism class of nonempty
objects, hence is equivalent to $[1]$. Localizing this category
either keeps it or also inverts the arrow from empty to nonempty,
giving a point. This proves the corrected collapse conclusion.
Conversely, the component-bijection axiom factors $\pi_0$ through
localization and separates all discrete cardinalities. These two
directions give precisely the corrected equivalences. The printed
proofs invoke product preservation at the step isolated above;
B.8 and B.16 show why the relative hypothesis cannot simply be
omitted from that step.
\end{proof}

\section{The fibre axioms with enough intrinsic fibrations}

\begin{claim}{C.1}{Proved}{Universal weakness and fibrewise invariance}
Suppose $\mathcal M$ has pullbacks, $W$ contains isomorphisms and
satisfies two-out-of-three, and every map factors as a map in $W$
followed by an intrinsic fibration. Then
\[
 W\cap\operatorname{Fib}_W=W^u.
\]
Moreover every $W$-map between intrinsic fibrations over $S$ remains
weak after arbitrary change of base. Thus XII axiom (b) follows.
\end{claim}
\begin{proof}
Let $u:X\to Y$ be weak over $S$, with both structural maps intrinsic,
and factor an arbitrary $T\to S$ as $T\xrightarrow i T_0
\xrightarrow p S$, with $i\in W$ and $p$ intrinsic. Pullback of $u$
over $T_0$ is weak by the definition of $p$. The comparison maps from
pullback over $T$ to pullback over $T_0$ are weak because they are
pullbacks of $i$ along base changes of $X\to S$ and $Y\to S$.
The commutative square and two-out-of-three show that the pullback of
$u$ over $T$ is weak. This proves (b).

Apply the assertion to the arrow $q\to1_S$ in the slice when
$q\in W\cap\operatorname{Fib}_W$; it says precisely that $q$ is
universally weak. Conversely, if $q:X\to S$ is universally weak,
then for a weak arrow $T'\to T$ over $S$ both vertical arrows in the
square with upper edge $X\times_ST'\to X\times_ST$ are weak.
Two-out-of-three makes that upper edge weak, proving intrinsic
membership; the identity base change gives $q\in W$.
\end{proof}

\begin{claim}{C.2}{Proved}{Lead 3 without functorial factorization}
Under C.1's hypotheses, for every full subcategory
\[
 \mathcal F_S:=\operatorname{Fib}_W S\ \subseteq\
 \mathcal H\ \subseteq\ \mathcal M/S
\]
the inclusions induce equivalences on ordinary localizations and on
$\infty$-categorical localizations at the restricted $W_S$.
Consequently factorization alone implies all three axioms (a), (b), (c).
No global terminal object, functorial factorization, retract closure,
or strong saturation of $W$ is required.
\end{claim}
\begin{proof}
We give all hypotheses of the localization theorem being used.
For an ordinary category $\mathcal A$, the nerve is an
$\infty$-category with weak equivalences and fibrations in the sense
of \citet[Definitions~7.4.6 and 7.4.12]{RESCisinski2019} if it has a
terminal object, the weak class contains identities and has
two-out-of-three, and a class $F$ containing identities and closed
under composition has the following properties. Pullbacks of
$F$-maps between fibrant objects along maps from fibrant objects
exist and stay in $F$; the same pullbacks preserve $F\cap W$;
every map with fibrant codomain factors as $W$ followed by $F$.
Here fibrant means that the terminal map belongs to $F$.
Cisinski's Theorem~7.5.18 then states that inclusion of the full
subcategory of fibrant objects induces an equivalence of
$\infty$-categorical localizations.

Take $\mathcal A=\mathcal H$, and let $F$ consist of its morphisms
whose underlying map in $\mathcal M$ is intrinsic. The object
$1_S$ is terminal in $\mathcal H$ and belongs to $\mathcal F_S$.
Identities and compositions are in $F$, by the corresponding
unconditional properties of intrinsic fibrations. The fibrant
objects of $\mathcal H$ are exactly $\mathcal F_S$.
Suppose $X\to Y$ is in $F$, with $X,Y$ fibrant over $S$, and
$Z\to Y$ has $Z$ fibrant. Its pullback $P\to Z$ is intrinsic,
so $P\to S$ is intrinsic by composition. Thus $P$ belongs to
$\mathcal F_S$, hence to $\mathcal H$. Since $\mathcal H$ is full,
this is also a pullback in $\mathcal H$. If $X\to Y$ is weak,
C.1 makes its pullback weak. Finally, for a map $B\to A$ in
$\mathcal H$ with $A$ fibrant, factor its underlying map as
$B\xrightarrow w A'\xrightarrow q A$, where $w\in W$ and $q$
is intrinsic. Composition with $A\to S$ makes $A'$ fibrant,
so the entire factorization lies in $\mathcal H$.

All of Cisinski's hypotheses have now been verified within
$\mathcal H$. Importantly, his definition does not require arbitrary
pullbacks in $\mathcal H$. The required pullbacks have fibrant
vertices and have just been shown to remain in the smaller
subcategory. The theorem gives $L(\mathcal F_S)\simeq L(\mathcal H)$.
Apply the same argument to $\mathcal M/S$ to obtain the composite
equivalence $L(\mathcal F_S)\simeq L(\mathcal M/S)$; the second
inclusion is consequently an equivalence as well. Taking homotopy
categories gives the ordinary Gabriel--Zisman localizations.
C.1 identifies $W_S$ with $W_S^u$ on $\mathcal F_S$, giving
axiom (a) with exactly its stated weak class.
\end{proof}

\begin{claim}{C.3}{Refuted}{Five objects suffice without factorization}
Axioms (a) and (b), even at every base and for strongly saturated $W$,
do not imply (c) or the intrinsic factorization hypothesis. A
counterexample has five objects and sixteen arrows. Thus the
six-object poset example is not minimal among finite categories.
\end{claim}
\begin{proof}
Let $\mathcal M$ have objects $0,A,B,C,S$. The object $0$ is initial,
with its unique arrows to the other objects. Besides these arrows
and the five identities, there are
\[
 \alpha,\beta:A\rightrightarrows B,\quad \gamma:A\to C,
 \quad b:B\to S,\quad c:C\to S,\quad x,y:A\rightrightarrows S,
\]
with the only nonidentity composition relations not involving $0$
given by $b\alpha=c\gamma=x$ and $b\beta=y$. Put
\[
 W=\{1_0,1_A,1_B,1_C,1_S,\alpha,\beta,\gamma\}.
\]
These formulas define the complete composition law; associativity
involves only identities or arrows from $0$ and follows directly.
The functor to the three-object chain $0<m<s$ that sends $A,B,C$
to $m$ and $S$ to $s$ has $W$ as the inverse image of identities.
It proves two-out-of-three, and also strong saturation: any arrow
inverted in $\mathcal M[W^{-1}]$ must map to an identity in this
chain, hence must already belong to $W$.

Each slice is a poset. Over $S$ its six objects are
$0_S,x,y,b,c,1_S$, with Hasse relations
\[
 0_S<x,y,\qquad x<b,c,\qquad y<b,\qquad b,c<1_S.
\]
The slice over $B$ is the diamond $0_B<\alpha,\beta<1_B$;
over $C$ it is the chain $0_C<\gamma<1_C$; over $A$ it is
$0_A<1_A$; over $0$ it is terminal. Each has all binary meets,
so every cospan of $\mathcal M$ has a pullback. In particular
$\alpha$ and $\beta$ pull back to $0$, while $b$ and $c$ pull
back to $A$ with projections $\alpha,\gamma$.

The intrinsic fibrant objects in the slices are exactly
\[
\begin{array}{c|l|l}
\text{base}&\text{slice objects}&\text{intrinsic fibrant objects}\\\hline
0&1_0&1_0\\
A&0_A,1_A&0_A,1_A\\
B&0_B,\alpha,\beta,1_B&0_B,1_B\\
C&0_C,\gamma,1_C&0_C,\gamma,1_C\\
S&0_S,x,y,b,c,1_S&0_S,b,1_S
\end{array}
\]
To check the only nontrivial row, pull back the three weak arrows
$x\to b$, $y\to b$, $x\to c$ in the slice over $S$.
The object $b$ preserves all three, producing $\alpha,\beta,1_A$.
The candidates $x,y,c$ fail, respectively, on $y\to b$,
$x\to b$, and $y\to b$, producing $0\to A$ in each case.
The bottom object and the identity always pass. The other rows
follow from the same diamond and chain tests. These tests include
all possible structure maps to the base, not only weak arrows
whose codomain is the base.

On the fibrant subcategory the weak arrows are only identities,
except at base $C$, where $\gamma$ is also weak. Its pullbacks
over $C$ are $1_0,1_A,\gamma$. This proves (b) at all bases.
For (a), only $B$ and $S$ need replacement. Over $B$ send
$\alpha,\beta$ to $1_B$ and fix $0_B,1_B$; the identity has a
weak natural transformation to this retraction. Over $S$ define
the two endofunctors by
\[
\begin{array}{c|rrrrrr}
z&0_S&x&y&b&c&1_S\\\hline
Rz&0_S&b&b&b&b&1_S\\
Tz&0_S&x&y&b&x&1_S .
\end{array}
\]
Both preserve the order and the weak arrows. There are natural
transformations $1\leftarrow T\to R$ with weak components, and
$R,T$ restrict to the identity on $\{0_S,b,1_S\}$. They exhibit
the localized inclusion of the fibrant subcategory as an
equivalence, at both the ordinary and the enhanced level.

The full intermediate subcategory on $\{0_S,b,c,1_S\}$ is a
diamond and has only identity weak arrows. Its fibrant subcategory
is the chain $\{0_S,b,1_S\}$, also with only identity weak arrows.
The inclusion fails essential surjectivity at $c$, so (c) fails.
Finally the only weak arrow with source $C$ is $1_C$, whereas
$c:C\to S$ is not intrinsic. It has no required factorization.
The program verifies the category laws, all sixty-six cospans and
their universal properties, every intrinsic test, every test in
(b), and both displayed deformation certificates. Its complete
output is \texttt{code/BC/verification\_output.json}.
\end{proof}

\begin{claim}{C.4}{Proved}{Why an object bound is not a finite category search bound}
There are infinitely many inequivalent finite categories with
pullbacks and one object. A finite terminating exhaustive search
through six objects therefore needs an additional arrow bound or
a structural reduction. Every finite category with pullbacks has
only monomorphisms, but this does not bound its number of arrows.
\end{claim}
\begin{proof}
The one-object categories of cyclic groups $C_m$, for arbitrary
positive integers $m$, have pullbacks: a cospan of invertible
arrows has a pullback obtained using either inverse. They are
pairwise inequivalent. To prove the structural assertion, every
slice of a finite category with pullbacks is finite and has
terminal object and binary products. For any objects $X,Y$ in
such a slice, the product powers $Y^n$ give
$|\operatorname{Hom}(X,Y^n)|=|\operatorname{Hom}(X,Y)|^n$.
There is a common finite bound on all hom-set cardinalities,
so $|\operatorname{Hom}(X,Y)|\le1$. Each slice is therefore a
preorder. If $fh=fk$, the two maps $h,k$ define parallel maps
in the slice over the codomain of $f$, and must coincide.
Thus $f$ is monic. In particular finite endomorphism monoids
are groups, but the preceding family shows that their orders
remain unbounded.
\end{proof}

\begin{claim}{C.5}{Proved}{A direct verification of the perfect path replacement}
For each small $S$ there is a strictly functorial coaugmented
replacement $R_S:\operatorname{Cat}/S\to\operatorname{Cat}/S$
whose values are $W_\infty$-smooth and proper and whose
coaugmentation has underlying map in $W_\infty$.
\end{claim}
\begin{proof}
Let $J$ be the integer fence, with an arrow from each even integer
to each adjacent odd integer. Let $\mathcal P(S)$ have the
eventually constant functors $J\to S$ and all natural transformations.
The left and right endpoint functors are $\sigma,\beta$. These
are defined on morphisms because naturality on identity tails
forces their components to be eventually constant too.

We check the incoming and outgoing lifting categories of
$(\sigma,\beta):\mathcal P(S)\to S\times S$ directly. For an
outgoing lift from $\gamma$ with prescribed endpoint arrows
$\ell\to\ell'$ and $r\to r'$, choose an odd left boundary $a$
and an even right boundary $b$, far enough out that $\gamma$
is already constant there. Use $\gamma$ on $[a+1,b]$, use
$\ell'$ on the left tail and $r'$ on the right tail, and use
the prescribed arrows at the joins $a+1\to a$ and $b\to b+1$.
The transformation from $\gamma$ is the identity on the middle
interval and the prescribed maps on the tails. In the full
stage of lifts stationary outside these boundaries this object
is initial: the map to another lift is its given transformation
on the middle interval and the identity on the two tails.
Naturality at the joins follows from the other lift's two
naturality squares; uniqueness follows from the identity middle
components of the universal lift. For incoming lifts choose
an even left boundary and an odd right boundary and reverse
this construction. It supplies a terminal object in each stage.
Every finite simplex of either lifting category occurs in one
stage, and enlarging the boundaries is filtered. Filtered
colimits of simplicial sets preserve weak equivalences.
Both full lifting categories are therefore weakly contractible.
The lifting-category criterion for smoothness and properness
proves the asserted two properties of the endpoint functor
\citep[Theorem~4.4.36(ii'), and its opposite]{RESCisinski2019}.

For $f:A\to S$ set
\[
 R_S(f)=\bigl(A\times_{S,\sigma}\mathcal P(S)
           \xrightarrow{\beta}S\bigr),
 \qquad \eta_f(a)=(a,\operatorname{const}_{f(a)}).
\]
The map to $A\times S$ is a base change of the endpoint functor,
and the projection $A\times S\to S$ is smooth and proper.
These properties are stable under composition and base change,
so $R_S(f)$ is perfect.

There is also a direct proof that the projection $R_S(f)\to A$
is weak. Restrict to a finite fence interval containing the
support of a path and extend constantly outside it. Deleting
its rightmost leaf, and then extending constantly from the
neighbor, gives a retraction of finite diagram categories.
The edge to the deleted leaf gives a natural transformation
in one of the two directions between this composite and the
identity. Repeat to contract the finite fence to its leftmost
vertex. Throughout, that vertex is $f(a)$ and the transformations
are the identity on $a$. Thus the finite-stage projection to
$A$ is a category-homotopy equivalence, with the constant-path
section. Passing to the filtered union proves the full
projection weak. Its section $\eta_f$ is weak by two-out-of-three.
For $u:A\to A'$ over $S$, send $(a,\gamma)$ to $(u(a),\gamma)$.
These formulas preserve composition and the coaugmentation
strictly. This supplies the promised functoriality and completes
an independent verification of the portion of BISITE B.19 used here.
\end{proof}

\begin{claim}{C.6}{Proved}{All intermediate perfect subcategories compute the same localization}
Let $W$ be any fundamental localizer. For either
$\mathcal P_S=\operatorname{Parf}_{W_\infty}S$ or
$\mathcal P_S=\operatorname{Parf}_W S$, and every full
$\mathcal P_S\subseteq\mathcal H\subseteq\operatorname{Cat}/S$,
both inclusions induce equivalences after localization at $W_S$,
ordinarily and as $\infty$-categories. This applies in particular
to all $W_n$ and every homology localizer.
\end{claim}
\begin{proof}
The minimal-localizer theorem gives $W_\infty\subseteq W$
\citep[Theorem~6.1.18]{RESCisinski2006}. The lifting tests also
give $\operatorname{Parf}_{W_\infty}\subseteq\operatorname{Parf}_W$.
C.5 therefore provides a functor $R_S$ taking every object into
either chosen full subcategory, and a natural transformation
$\eta:1\to R_S$ with components in $W_S$. Naturality and
two-out-of-three show that $R_S$ preserves $W_S$.
For an object of $\mathcal H$, both ends of $\eta$ lie in
$\mathcal H$, so the transformation is internal to it by
fullness. The same holds in $\mathcal P_S$.
After localization, these transformations are natural
equivalences and make $R_S$ an inverse to each inclusion.
No assertion that all perfect maps are intrinsic for $W$
enters this argument. In particular the failure of that
intrinsic assertion for homology localizers does not obstruct
the present conclusion.
\end{proof}

\begin{claim}{C.7}{Proved}{The precise conclusion at XII ms.~240--241}
The left factorization hypothesis at XII ms.~240 implies the
three fibre axioms as printed at ms.~215--218. For arbitrary
$g:T\to S$, pullback on intrinsic fibrant objects defines a
derived reindexing functor on the localized slices. For a
strict pullback $Y=X_1\times_S X_2$ with $X_1\to S$ intrinsic,
both displayed relative homotopy-fibre identities at ms.~241
hold. These assertions do not require every terminal map to
be intrinsic. The weaker existence condition in footnote~128,
that intrinsic maps be pullbackable, suffices.
\end{claim}
\begin{proof}
C.1 and C.2 prove the three axioms. Pullback along $g$ takes
intrinsic fibrations to intrinsic fibrations and preserves
their weak arrows by C.1, so it descends to their localizations.
C.2 identifies these with the localized full slices. The
canonical isomorphisms between iterated pullbacks provide the
coherence for successive base changes.

The first identity is computed by $X_1\to S$ itself, which
is already fibrant. For the second, choose a factorization
$X_2\xrightarrow w\bar X_2\to S$ with $w\in W$ and the
second map intrinsic. The comparison
$X_1\times_SX_2\to X_1\times_S\bar X_2$ belongs to $W$
by the definition of the intrinsic map $X_1\to S$.
It therefore identifies the raw pullback with the replacement
calculation in the localized slice over $X_1$, exactly as
required. This also proves independence of replacement
through the equivalences in C.2.

Inspecting these proofs, every pullback used is along an
intrinsic map, or along a composite or a base change of one.
Those maps need only be pullbackable; arbitrary pullbacks
elsewhere in $\mathcal M$ are unnecessary. The terminal
object needed in the localization argument is $1_S$ in the
slice. Thus the caution in footnote~130 about all objects
being globally fibrant is unnecessary for the stated
relative assertions. The word ``cofibrations'' in
footnote~129 has the opposite variance from the displayed
factorization and should read ``fibrations.'' The two
mislabelled arrows at ms.~241 are $\bar p_2$, not $p_2$
and $\bar q_2$, respectively, as follows from their domains
and codomains. These are mathematical typing corrections;
no additional palaeographical claim is made.
\end{proof}

\begin{claim}{C.8}{Open}{The remaining minimum problem}
The least number of objects of a finite category with pullbacks
admitting (a) and (b) but not (c) has not been determined here.
C.3 proves an upper bound of five. The unbounded-arrow
classification required to exclude all smaller finite
categories has not been supplied. No exhaustive enumeration
through six objects is claimed.
\end{claim}
